\section{Research Gap and Related Work}
\label{sec:related}

\paragraph{SOC assessment and performance measurement.}
SOC-CMM scores capability maturity across business, people, process, technology
and services domains through a self-assessment tool~\citep{vanos2016soccmm}.
Agyepong et al.\ weight analyst-performance indicators using a Delphi panel and
the analytic hierarchy process~\citep{agyepong2023analyst}; Forsberg and Frantti
derive technical SOC metrics through design science~\citep{forsberg2023metrics}.
Surveys and qualitative studies characterise SOC building blocks, challenges and
human factors~\citep{vielberth2020soc,kokulu2019soc,sundaramurthy2015burnout,sundaramurthy2016contradictions}.
These instruments measure maturity, define metrics or explain practice; they do
not compute findings from the operational records a monitored SOC submits. SAT-SA
addresses that narrower gap: given submitted records, which supervisory concerns
are supported by which records?

\paragraph{Alert prioritisation and alert fatigue.}
A large literature ranks alerts or incidents for SOC analysts. Jalalvand et
al.\ review 89 prioritisation papers through a human--AI teaming lens~\citep{jalalvand2024prioritisation};
Tariq et al.\ survey alert-fatigue mitigations~\citep{tariq2025alertfatigue}.
Systems include provenance-based triage (NoDoze~\citep{hassan2019nodoze}), guided
response trained on customer triage labels~\citep{freitas2024guide}, large-scale
incident ranking~\citep{freitas2026aip} and learning to defer to analysts~\citep{jalalvand2025l2d}.
SAT-SA differs in unit and user: it ranks \emph{monitored organisations} for a
\emph{supervisor} from rule-based findings about how work was handled, not
individual alerts for an analyst. It is not a competitor to these systems and is
not compared with them.

\paragraph{Process analytics and absence detection.}
Process mining discovers and checks processes from event logs~\citep{vanderaalst2016pm};
conformance checking relates recorded behaviour to a model and exposes skipped or
out-of-order activities~\citep{carmona2018conformance}. It has been applied to
security audit trails~\citep{vanderaalst2005security} and to assessing IT incident
management against ITIL~\citep{ferreira2008itil}. SAT-SA's execution-gap and
``negative-space'' workers are rule-encoded conformance checks (for example, an
acknowledged alert whose case has no investigation step); we treat them as an
adaptation of this prior art, not a new technique.

\paragraph{Peer comparison and robust deviation.}
Peer group analysis flags entities that diverge from similar entities over
time~\citep{weston2008pga}, in the broader statistical fraud-detection
tradition~\citep{bolton2002fraud}. SAT-SA's peer worker compares an entity with a
declared cohort using median and median absolute deviation (MAD)~\citep{leys2013mad},
with a minimum cohort size and tenant isolation.

\paragraph{Supervisory technology.}
Financial supervisors increasingly use technology to analyse regulated
entities' data (``suptech'')~\citep{broeders2018suptech}; human supervisory
judgement remains indispensable there, as in SAT-SA's design.

\paragraph{Integrity and provenance of records.}
Forward-secure and tamper-evident logs~\citep{schneier1999logs,crosby2009logging}
detect alteration of recorded history; W3C PROV provides a provenance
vocabulary~\citep{w3c2013provdm}. SAT-SA composes standard primitives---SHA3-256~\citep{nist2015fips202}
and ML-DSA-65~\citep{nist2024fips204}---to bind a canonical supervisory record and the
authorised decision; it does not propose cryptography. Automation bias in human
use of decision support~\citep{parasuraman2010automation} motivates keeping the
decision with the supervisor.

\paragraph{Gap, stated conservatively.}
We did not find, in this focused review, a system that computes evidence-cited
supervisory findings from SOC submission records for an external supervisor and
binds the supervisor's decision to that analysed state. This is an absence in a
non-systematic search, not evidence of novelty; the framing is therefore offered
as a low-confidence candidate contribution.
